\documentclass[aspectratio=1610]{beamer}
\usepackage{tikz}
\usepackage{booktabs}
\usepackage{graphicx}
\usepackage{amsmath}
\usepackage{appendixnumberbeamer}
\usetheme{Madrid}
\usecolortheme{default}

% Figure path macros (relative to presentation directory)
\newcommand{\kinetixfig}[1]{../report/sources/bt-kinetix/bt-kinetix/bt_experiments/final_plots/#1}
\newcommand{\liberofig}[1]{../report/sources/bt-libero/final_plots/#1}
\newcommand{\liberoablationfig}[1]{../report/sources/bt-libero/final_plots_vlash_d0_d4_d8_d16_over_steps/#1}

\title[Asynchronous VLA Inference]
{Bachelor Thesis}
\subtitle{Understanding Asynchronous Inference Methods\\for Vision-Language-Action Models}
\author[]
{Ayoub Agouzoul\inst{1,2}}

\institute[]
{
    \inst{1}%
    \'Ecole polytechnique\\
    \inst{2}%
    ETH Z\"urich\\[0.5em]
    \textit{Supervised by Professor Wenqi Jiang}
}

\date[March 27th, 2026]
{January---March 2026}

% TODO: add supervisor + referent instructor

\logo{} % empty initially; activated after title slide
\titlegraphic{%
  \raisebox{-0.5\height}{\includegraphics[width=0.12\textwidth]{logo-EP-vertical.png}}%
  \hspace{2em}%
  \raisebox{-0.5\height}{\includegraphics[width=0.12\textwidth]{systemslogo-3colour-blacktext.pdf}}%
}

%------------------------------------------------------------
% Table of contents at each section
\AtBeginSection[]
{
  \begin{frame}
    \frametitle{Outline}
    \tableofcontents[currentsection]
  \end{frame}
}
%------------------------------------------------------------

\begin{document}

\frame{\titlepage}

% Activate logo for all subsequent slides
\logo{\includegraphics[width=0.1\textwidth]{logo-EP-vertical.png}}

%------------------------------------------------------------
% OUTLINE
%------------------------------------------------------------
\begin{frame}
\frametitle{Outline}
\tableofcontents
\end{frame}

%============================================================
\section{Motivation}
%============================================================

%------------------------------------------------------------
\begin{frame}
\frametitle{VLA Models Promise Generalist Robot Control}

\textbf{Vision-Language-Action (VLA) models:}
\begin{itemize}
    \item Map camera images + language instructions $\rightarrow$ robot actions
    \item Leverage pretrained vision-language representations
    \item State-of-the-art: $\pi_0$, $\pi_{0.5}$, SmolVLA
\end{itemize}

\vspace{0.5em}
\textbf{Action chunking} for efficient deployment:
\begin{itemize}
    \item Predict $H$ future actions in one forward pass (via flow matching)
    \item Execute chunk while computing the next one (asynchronous)
\end{itemize}

\vspace{0.5em}
\begin{alertblock}{Key problem}
Large model $\rightarrow$ high inference latency $\rightarrow$ by the time actions are ready, the robot's state has changed
\end{alertblock}

\end{frame}

%------------------------------------------------------------
\begin{frame}
\frametitle{Asynchronous Inference Creates Observation Staleness}

\textbf{Inference delay} $d = \lfloor \delta / \Delta t \rfloor$:
\begin{itemize}
    \item $\delta$: wall-clock inference time
    \item $\Delta t$: control period (e.g., 20\,ms at 50\,Hz)
    \item First $d$ actions in each new chunk are \alert{already obsolete}
\end{itemize}

\vspace{0.5em}
\textbf{Naive async execution causes:}
\begin{itemize}
    \item \textbf{Chunk boundary discontinuity} --- abrupt transitions, jitter
    \item \textbf{Observation staleness} --- actions conditioned on outdated state
    \item \textbf{Open-loop degradation} --- no correction during chunk lifetime
\end{itemize}

\vspace{0.5em}
\begin{block}{How severe?}
Cloud inference over 4G: $d = 13$ steps; edge GPU on-device: $d = 2$
\end{block}

\end{frame}

%------------------------------------------------------------
\begin{frame}
\frametitle{Four Methods Address Staleness --- Which Works Best?}

\begin{enumerate}
    \item \textbf{IT-RTC:} Inpaints fresh actions at inference time (training-free)
    \item \textbf{TT-RTC:} Simulates delays during training (zero inference overhead)
    \item \textbf{VLASH:} Conditions on estimated future robot state
    \item \textbf{A2C2:} Lightweight correction head for per-step residual fixes
\end{enumerate}

\vspace{0.5em}
\textbf{Problem:} Developed independently, incompatible codebases and protocols

\vspace{0.5em}
\begin{block}{Our contribution}
\begin{itemize}
    \item \textbf{Two unified codebases} for fair side-by-side comparison
    \item \textbf{Systematic evaluation} on Kinetix and LIBERO (delays up to $d = 20$)
    \item \textbf{Cost analysis and practical guidelines} for choosing a method
\end{itemize}
\end{block}

\end{frame}

%============================================================
\section{Background}
%============================================================

%------------------------------------------------------------
\begin{frame}
\frametitle{How Do VLA Models Generate Actions?}

\textbf{Architecture (e.g., SmolVLA, $\pi_0$):}
\begin{itemize}
    \item \textbf{VLM backbone:} Processes images + language $\rightarrow$ features
    \item \textbf{Action expert:} Generates continuous action chunk via flow matching
\end{itemize}

\vspace{0.3em}
\textbf{Flow matching} for action generation:
\begin{itemize}
    \item Start from noise $\epsilon \sim \mathcal{N}(0, I)$
    \item Iteratively denoise through $n$ Euler steps (typically 5--10)
    \item $A^{\tau + 1/n} = A^{\tau} + \frac{1}{n} v_\theta(A^{\tau}, o, \tau)$
\end{itemize}

\vspace{0.3em}
\textbf{Action chunking:}
\begin{itemize}
    \item Predict $H$ future actions: $A_t = (a_t, \ldots, a_{t+H-1})$
    \item Execute $s \leq H$ actions, then replan with fresh observation
\end{itemize}

\end{frame}

%------------------------------------------------------------
\begin{frame}
\frametitle{IT-RTC and TT-RTC Fix Chunk Boundary Discontinuity}

\begin{columns}
\column{0.5\textwidth}
\textbf{IT-RTC} (inference-time):
\begin{itemize}
    \item Treats committed prefix as constraint
    \item Soft mask + pseudoinverse VJP guidance during denoising
    \item \alert{No retraining needed}
    \item Overhead: $3\times$ FLOPs on denoising
\end{itemize}

\column{0.5\textwidth}
\textbf{TT-RTC} (training-time):
\begin{itemize}
    \item Simulates random delays $d \sim p(d)$ during training
    \item Per-token flow timesteps: prefix at $\tau{=}1$, postfix sampled
    \item Loss computed only on postfix
    \item \alert{Zero inference overhead}
\end{itemize}
\end{columns}

\vspace{0.5em}
\begin{block}{Key difference}
IT-RTC approximates via linearized VJP guidance; TT-RTC learns the prefix-conditioned distribution directly
\end{block}

\end{frame}

%------------------------------------------------------------
\begin{frame}
\frametitle{VLASH and A2C2 Target Different Failure Modes}

\begin{columns}
\column{0.5\textwidth}
\textbf{VLASH} (future-state-aware):
\begin{itemize}
    \item Estimates robot state at execution time:
    $\hat{s}_{t+\Delta} = s_t + \sum_k a_{t+k}^{\text{prev}}$
    \item Conditions policy on future state
    \item Temporal-offset augmentation during training
    \item Needs dynamics model or simulator
\end{itemize}

\column{0.5\textwidth}
\textbf{A2C2} (residual correction):
\begin{itemize}
    \item Lightweight correction head runs every control step
    \item $a^{\text{exec}} = a^{\text{base}} + \Delta a$
    \item Uses \emph{latest} observation
    \item \alert{Plug-in:} no base policy retraining
    \item Can combine with other methods
\end{itemize}
\end{columns}

\vspace{0.5em}
\begin{block}{Each method targets a different failure mode}
IT-RTC \& TT-RTC $\rightarrow$ chunk discontinuity, \quad VLASH $\rightarrow$ staleness, \quad A2C2 $\rightarrow$ open-loop degradation
\end{block}

\end{frame}

%============================================================
\section{Experimental Setup}
%============================================================

%------------------------------------------------------------
\begin{frame}
\frametitle{Unified Codebases Enable Fair Comparison}

\textbf{Problem:} Each method in separate codebase with different dependencies

\vspace{0.5em}
\begin{columns}
\column{0.5\textwidth}
\textbf{bt-kinetix} (JAX):
\begin{itemize}
    \item All 4 methods + naive baseline
    \item MLP-Mixer backbone ($\sim$310K params)
    \item Chunk sizes $H \in \{16, 30\}$
    \item 2048 parallel rollouts per config
\end{itemize}

\column{0.5\textwidth}
\textbf{bt-libero} (PyTorch / LeRobot):
\begin{itemize}
    \item SmolVLA base policy (604M params)
    \item All 4 methods + naive baseline
    \item Unified evaluation harness
    \item 4 LIBERO task suites (40 tasks)
\end{itemize}
\end{columns}

\vspace{0.5em}
\begin{block}{Key benefit}
Same base policy, data, and evaluation protocol $\rightarrow$ performance differences attributable to method, not implementation
\end{block}

\end{frame}

%------------------------------------------------------------
\begin{frame}
\frametitle{Two Benchmarks Cover Different Scales and Tasks}

\begin{columns}
\column{0.5\textwidth}
\textbf{Kinetix} (2D physics):
\begin{itemize}
    \item 10 dynamic environments
    \item Throwing, catching, balancing, locomotion
    \item Symbolic observations ($\text{dim} = 679$)
    \item Fast JAX simulation on GPU
    \item Delays $d = 0, \ldots, 15$
\end{itemize}

\column{0.5\textwidth}
\textbf{LIBERO} (manipulation):
\begin{itemize}
    \item 40 language-conditioned tasks
    \item 7-DoF Franka Panda arm
    \item RGB images + proprioception
    \item SmolVLA (production-scale VLA)
    \item Delays $d = 0, \ldots, 20$
\end{itemize}
\end{columns}

\vspace{0.5em}
\textbf{Research questions:}
\begin{enumerate}
    \item[\textbf{RQ1}] How do methods compare as delay $d$ increases?
    \item[\textbf{RQ2}] Do training-based methods generalize beyond $d_{\max}$?
    \item[\textbf{RQ3}] What are the training and inference costs of each method?
\end{enumerate}

\end{frame}

%------------------------------------------------------------
\begin{frame}
\frametitle{Real-World Delays Motivate Extended Evaluation Range}

\begin{table}
\centering
\small
\begin{tabular}{l l r r}
\toprule
\textbf{Scenario} & \textbf{Configuration} & \textbf{Latency $\delta$} & \textbf{Delay $d$} \\
\midrule
Consumer GPU (local) & RTX 4090 & 31\,ms & 1 \\
Edge GPU (on-device) & Jetson Thor & 53\,ms & 2 \\
Edge server (4G) & B100 + 4G & 73\,ms & 3 \\
Cloud server (slow) & B100 + 4G + cloud & 273\,ms & \alert{13} \\
\bottomrule
\end{tabular}
\end{table}

\small Based on $\pi_0$ latency estimates (Jiang et~al.), 50\,Hz controller

\vspace{0.5em}
\begin{alertblock}{Prior work only evaluated $d \leq 4$}
We extend to $d = 15$ (Kinetix) and $d = 20$ (LIBERO) to cover realistic deployment scenarios
\end{alertblock}

\end{frame}

%============================================================
\section{Results}
%============================================================

%------------------------------------------------------------
\begin{frame}
\frametitle{A2C2 Leads; TT-RTC Surpasses IT-RTC ($H = 16$)}

\begin{center}
    \includegraphics[width=0.95\textwidth,height=0.68\textheight,keepaspectratio]{\kinetixfig{c16/combined_2x2.pdf}}
\end{center}

\vspace{-0.5em}
\small
\textbf{Kinetix, 10 envs, 2048 rollouts each.} A2C2 (pink) leads, above 90\% even at $d = 8$. VLASH outperforms RTC methods. TT-RTC $d_{\max}{=}8$ surpasses IT-RTC at higher delays. Naive (red) drops below 40\%.

\end{frame}

%------------------------------------------------------------
\begin{frame}
\frametitle{IT-RTC Degrades at High Delays ($H = 30$)}

\begin{center}
    \includegraphics[width=0.95\textwidth,height=0.68\textheight,keepaspectratio]{\kinetixfig{c30/combined_2x2.pdf}}
\end{center}

\vspace{-0.5em}
\small
\textbf{Extended to $d = 15$.} A2C2 still leads. \alert{IT-RTC falls close to naive} --- VJP approximation degrades as prefix-to-postfix ratio grows. TT-RTC remains robust. VLASH $d_{\max}{=}15$ hurts low-delay performance ($\sim$10\% drop).

\end{frame}

%------------------------------------------------------------
\begin{frame}
\frametitle{TT-RTC Generalizes Stably; VLASH Trades Low for High Delay}

\begin{center}
    \includegraphics[width=0.95\textwidth,height=0.68\textheight,keepaspectratio]{\kinetixfig{ttrtc_vlash_c30_d5_vs_d9_vs_d16/combined_2x2.pdf}}
\end{center}

\vspace{-0.5em}
\small
\textbf{$d_{\max}$ ablation ($H = 30$).}
TT-RTC variants evolve very similarly regardless of $d_{\max}$.
VLASH shows a clear \alert{low-delay vs.\ high-delay trade-off}: wider $d_{\max}$ helps at high $d$ but hurts near $d = 0$.

\end{frame}

%------------------------------------------------------------
\begin{frame}
\frametitle{Results Extend to Production VLA Models on LIBERO}

\vspace{-0.5em}
\begin{center}
    \includegraphics[width=0.70\textwidth,height=0.55\textheight,keepaspectratio]{\liberofig{success_vs_delay_action_50_averaged.pdf}}
\end{center}

\vspace{-0.5em}
\small
SmolVLA (604M params) on 40 LIBERO tasks, $s = 50$, delays $d \in \{0, \ldots, 20\}$

\end{frame}

%------------------------------------------------------------
\begin{frame}
\frametitle{VLA-Scale Cost Analysis Reveals Surprising Trade-offs}

\begin{table}
\centering
\footnotesize
\begin{tabular}{l r r}
\toprule
\textbf{Method} & \textbf{Per-chunk GFLOPs} & \textbf{Ratio} \\
\midrule
Naive / TT-RTC / VLASH & 551.0 & $1.00\times$ \\
IT-RTC & 663.8 & $1.20\times$ \\
A2C2 ($s{=}50$) & \alert{961.7} & \alert{$1.75\times$} \\
\bottomrule
\end{tabular}
\end{table}

\vspace{0.3em}
\textbf{Why IT-RTC is only $1.2\times$ (not $3\times$)?}
\begin{itemize}
    \item Vision encoder prefill = \textbf{89.8\%} of chunk cost (494.6 / 551.0\,GFLOPs)
    \item IT-RTC's $3\times$ overhead applies only to the denoising loop (10\% of cost)
\end{itemize}

\vspace{0.3em}
\textbf{Why A2C2 is the most expensive?}
\begin{itemize}
    \item 50 ResidualTransformer corrections $\times$ 8.21\,GFLOPs each = 410.7\,GFLOPs
\end{itemize}

\vspace{0.3em}
\begin{block}{Latency still matters}
IT-RTC triples the \emph{sequential} denoising bottleneck $\rightarrow$ avoid for real-time
\end{block}

\end{frame}
%============================================================
\section{Conclusion}
%============================================================

%------------------------------------------------------------
\begin{frame}
\frametitle{Summary: No Single Winner, But Clear Trade-offs}

\textbf{Key findings:}
\begin{itemize}
    \item All four methods substantially outperform naive async execution
    \item A2C2's per-step correction most effective ($>$90\% at $d{=}0...8$ on Kinetix)
    \item IT-RTC degrades at high delays and long chunks ($H = 30$)
    \item TT-RTC is the most \alert{robust}: stable across $d_{\max}$, zero overhead
    \item VLASH has a clear low-delay vs.\ high-delay trade-off with $d_{\max}$
\end{itemize}

\vspace{0.5em}
\textbf{Practical guidelines:}
\begin{itemize}
    \item No retraining possible? $\rightarrow$ \textbf{IT-RTC} (drop-in, but degrades at high $d$)
    \item Fine-tuning feasible? $\rightarrow$ \textbf{TT-RTC} (best robustness--cost trade-off)
    \item Dynamics model available? $\rightarrow$ \textbf{VLASH} (tune $d_{\max}$ carefully)
    \item Want per-step corrections? $\rightarrow$ \textbf{A2C2} (plug-in, combinable)
\end{itemize}

\end{frame}

%------------------------------------------------------------
\begin{frame}
\frametitle{Limitations}

\begin{itemize}
    \item VLASH on LIBERO uses ground-truth future states (simulator), not available in real deployment
    \item Hyperparameters from original papers, not tuned per benchmark
    \item Inference delay $d$ fixed per episode; real deployments have variable latency
    \item Limited GPU budget: not all method combinations and chunk sizes explored
    \item Walltime measurements on Kinetix not representative of real-world robot deployment
\end{itemize}

\end{frame}

%------------------------------------------------------------
\begin{frame}
\frametitle{Thank You --- Questions?}

\begin{center}
\Large
\textbf{Questions and Discussion}
\end{center}

\vspace{1em}
\textbf{Main takeaways:}
\begin{enumerate}
    \item Observation staleness is a critical VLA deployment bottleneck
    \item No single method wins everywhere --- the right choice depends on deployment constraints
    \item Unified evaluation reveals trade-offs invisible when methods are compared separately
\end{enumerate}

\vspace{1em}
\begin{block}{Code available}
\url{https://github.com/TheAyos/async-vla-inference}
\end{block}

\begin{block}{Acknowledgments}
Prof.\ Wenqi Jiang, Prof.\ Sergio Mover, ETH Systems Group, \'Ecole polytechnique
\end{block}

\end{frame}

%============================================================
% APPENDIX (backup slides for questions)
%============================================================
\appendix

%------------------------------------------------------------
\begin{frame}
\frametitle{Appendix: Deployment Scenario Delay Mapping}

\begin{table}
\centering
\small
\begin{tabular}{l l r r}
\toprule
\textbf{Scenario} & \textbf{Configuration} & \textbf{Latency $\delta$} & \textbf{Delay $d$} \\
\midrule
Datacenter GPU (local) & B100, no network & 3.2\,ms & 0 \\
Consumer GPU (local) & RTX 4090, no network & 31.1\,ms & 1 \\
Edge GPU (on-device) & Jetson Thor, no network & 52.6\,ms & 2 \\
Edge server (WiFi\,7) & B100 + WiFi\,7 & 8.4\,ms & 0 \\
Edge server (5G) & B100 + 5G & 27.8\,ms & 1 \\
Edge server (4G) & B100 + 4G & 73.0\,ms & 3 \\
Cloud server (fast) & B100 + wired + fast cloud & 23.4\,ms & 1 \\
Cloud server (slow) & B100 + 4G + slow cloud & 273.4\,ms & 13 \\
\bottomrule
\end{tabular}
\end{table}

\small Based on $\pi_0$ latency estimates from Jiang et~al., assuming 50\,Hz control frequency ($\Delta t = 20$\,ms).

\end{frame}


%------------------------------------------------------------
\begin{frame}
\frametitle{Appendix: Baseline Sanity Check --- $d = 0$ Performance}

\vspace{-0.5em}
\begin{center}
    \includegraphics[width=\textwidth,height=0.36\textheight,keepaspectratio]{\kinetixfig{c16/baseline_comparison.pdf}}

    \vspace{0.1em}

    \includegraphics[width=\textwidth,height=0.36\textheight,keepaspectratio]{\kinetixfig{c30/baseline_comparison.pdf}}
\end{center}

\vspace{-0.5em}
\small All methods achieve comparable $d = 0$ performance. Minor variations for VLASH with high $d_{\max}$.

\end{frame}

%------------------------------------------------------------
\begin{frame}
\frametitle{Appendix: Cross-Chunk-Size Comparison (Naive, IT-RTC, A2C2)}

\begin{center}
    \includegraphics[width=0.95\textwidth,height=0.72\textheight,keepaspectratio]{\kinetixfig{c16_vs_c30_naive_itrtc_residual/combined_2x2.pdf}}
\end{center}

\small $H = 16$ consistently outperforms $H = 30$ --- longer chunks impose intrinsic prediction difficulty.

\end{frame}

%------------------------------------------------------------
\begin{frame}
\frametitle{Appendix: Cross-Chunk-Size Comparison (TT-RTC)}

\begin{center}
    \includegraphics[width=0.95\textwidth,height=0.72\textheight,keepaspectratio]{\kinetixfig{c16_vs_c30_ttrtc/combined_2x2.pdf}}
\end{center}

\small TT-RTC's delay robustness transfers across chunk sizes when $d_{\max}$ is set appropriately.

\end{frame}

%------------------------------------------------------------
\begin{frame}
\frametitle{Appendix: Cross-Chunk-Size Comparison (VLASH)}

\begin{center}
    \includegraphics[width=0.95\textwidth,height=0.72\textheight,keepaspectratio]{\kinetixfig{c16_vs_c30_vlash/combined_2x2.pdf}}
\end{center}

\small VLASH's state rollforward accuracy is more sensitive to chunk size than TT-RTC's prefix conditioning.

\end{frame}

%------------------------------------------------------------
\begin{frame}
\frametitle{Appendix: Episode Efficiency ($H = 16$)}

\begin{center}
    \includegraphics[width=0.95\textwidth,height=0.72\textheight,keepaspectratio]{\kinetixfig{c16/episode_lengths_vs_delay_and_horizon_2x2.pdf}}
\end{center}

\small A2C2 achieves the lowest time-to-solve among solved episodes. Naive requires the most steps.

\end{frame}

%------------------------------------------------------------
\begin{frame}
\frametitle{Appendix: LIBERO Per-Suite Success Rate Breakdown}

\begin{center}
    \includegraphics[width=\textwidth,height=0.78\textheight,keepaspectratio]{\liberofig{success_vs_delay_action_50.pdf}}
\end{center}

\end{frame}

%------------------------------------------------------------
\begin{frame}
\frametitle{Appendix: LIBERO Per-Suite Episode Length}

\begin{center}
    \includegraphics[width=\textwidth,height=0.78\textheight,keepaspectratio]{\liberofig{solved_ep_length_vs_delay_action_50.pdf}}
\end{center}

\end{frame}

%------------------------------------------------------------
\begin{frame}
\frametitle{Appendix: VLASH Training Ablation on LIBERO}

\begin{center}
    \includegraphics[width=0.75\textwidth,height=0.68\textheight,keepaspectratio]{\liberoablationfig{success_vs_delay_action_30_averaged.pdf}}
\end{center}

\small VLASH variants with $d_{\max} \in \{0, 4, 8, 16\}$ at multiple training checkpoints (2k--5k steps), execution horizon $s = 30$.

\end{frame}

%------------------------------------------------------------
\begin{frame}
\frametitle{Appendix: Kinetix Inference FLOPs per Method}

\begin{table}
\centering
\small
\begin{tabular}{l l r r}
\toprule
\textbf{Method} & \textbf{Per-chunk FLOPs} & \textbf{MFLOPs} & \textbf{Ratio} \\
\midrule
Naive     & $N_s \cdot F_{\text{fwd}}$ & 235 & $1.0\times$ \\
IT-RTC    & $3 N_s \cdot F_{\text{fwd}}$ & 707 & $3.0\times$ \\
TT-RTC    & $N_s \cdot F_{\text{fwd}}$ & 235 & $1.0\times$ \\
VLASH     & $N_s \cdot F_{\text{fwd}}$ (+ env sim) & 235 & $1.0\times$ \\
A2C2      & $N_s \cdot F_{\text{fwd}} + s \cdot F_{\text{res}}$ & $235 + 0.62s$ & $\approx 1.02\times$ \\
\bottomrule
\end{tabular}
\end{table}

\small $C{=}16$, $N_s{=}5$ Euler steps, $F_{\text{fwd}} \approx 47.1$\,MFLOPs, $F_{\text{res}} \approx 0.62$\,MFLOPs.

On Kinetix (no vision encoder), IT-RTC's $3\times$ overhead is a true $3\times$. At VLA scale, the expensive vision prefill amortizes it to only $1.2\times$.

\end{frame}

%------------------------------------------------------------
\begin{frame}
\frametitle{Appendix: SmolVLA FLOPs Breakdown (LIBERO)}

\begin{table}
\centering
\small
\begin{tabular}{l r}
\toprule
\textbf{Component} & \textbf{FLOPs} \\
\midrule
SmolVLA prefill ($F_{\text{pre}}$) & 494.6\,GFLOPs \\
SmolVLA denoising step ($F_{\text{ds}}$) & 5.64\,GFLOPs \\
SmolVLA full chunk ($F_{\text{pre}} + 10 F_{\text{ds}}$) & 551.0\,GFLOPs \\
\midrule
ResidualTransformer forward ($F_{\text{res}}^L$) & 8.21\,GFLOPs \\
\midrule
SmolVLA total parameters & 604.9\,M \\
ResidualTransformer (total / trainable) & 32.0\,M / 20.8\,M \\
\bottomrule
\end{tabular}
\end{table}

\textbf{Key insight:} Vision encoder prefill = \textbf{89.8\%} of total chunk cost. The denoising loop (where IT-RTC's overhead applies) is only $\sim$10\%.

\end{frame}

%------------------------------------------------------------
\begin{frame}
\frametitle{Appendix: Measured LIBERO Inference Latency}

\begin{table}
\centering
\small
\begin{tabular}{l r r r r}
\toprule
\textbf{Method} & \textbf{Mean (ms)} & \textbf{P95 (ms)} & \textbf{Hz} & \textbf{Ratio} \\
\midrule
Naive & 405.2 & 407.5 & 2.47 & $1.00\times$ \\
TT-RTC & 402.7 & 404.9 & 2.48 & $0.99\times$ \\
IT-RTC & 469.7 & 488.6 & 2.13 & $1.16\times$ \\
A2C2 (serial) & 412.4 & 417.2 & 2.42 & $1.02\times$ \\
A2C2 residual only & 7.27 & 7.35 & 137.5 & $0.02\times$ \\
\bottomrule
\end{tabular}
\end{table}

\small RTX 3090, SmolVLA, batch size 1, 100 timed iterations

\vspace{0.3em}
\textbf{Key takeaways:}
\begin{itemize}
    \item TT-RTC indistinguishable from naive (402.7 vs.\ 405.2\,ms)
    \item IT-RTC +15.9\% latency (2.47 $\rightarrow$ 2.13\,Hz) --- slowest method
    \item A2C2 residual head cheap in isolation (7.27\,ms), serial path close to naive
\end{itemize}

\end{frame}

%------------------------------------------------------------
\begin{frame}
\frametitle{Appendix: Break-Even Analysis (TT-RTC vs.\ IT-RTC)}

\textbf{When does TT-RTC's training investment pay off?}
\[
N^* = \frac{C_{\text{train}}^{\text{TT-RTC}}}{\Delta C_{\text{chunk}}} \quad\text{where}\quad \Delta C_{\text{chunk}} = C_{\text{chunk}}^{\text{IT-RTC}} - C_{\text{chunk}}^{\text{TT-RTC}}
\]

\vspace{0.3em}
\textbf{LIBERO (SmolVLA):}
\begin{itemize}
    \item $C_{\text{train}}^{\text{TT-RTC}} \approx 1{,}310$\,PFLOPs (dominated by vision encoder forward)
    \item $\Delta C_{\text{chunk}} = 20 F_{\text{ds}} = 112.8$\,GFLOPs (only denoising saved)
    \item $N^* \approx 11.6$\,M chunks ($\approx 2.9$\,M episodes)
\end{itemize}

\vspace{0.3em}
\textbf{In pure FLOPs:} IT-RTC is surprisingly efficient at VLA scale ($1.2\times$ overhead)

\textbf{In latency:} IT-RTC triples the sequential denoising bottleneck $\rightarrow$ \textbf{avoid for real-time}

\end{frame}

%------------------------------------------------------------
\begin{frame}
\frametitle{Appendix: Flow Matching for Action Generation}

\textbf{Interpolant:} $A^\tau = (1 - \tau)\epsilon + \tau A_0$, \quad $\epsilon \sim \mathcal{N}(0, I)$

\vspace{0.3em}
\textbf{Training loss:}
\[
\mathcal{L}_{\text{FM}} = \mathbb{E}_{\tau, \epsilon}\left[\|v_\theta(A^\tau, o, \tau) - (A_0 - \epsilon)\|^2\right]
\]

\vspace{0.3em}
\textbf{Inference} (Euler integration from noise):
\[
A^{\tau + 1/n} = A^\tau + \frac{1}{n}v_\theta(A^\tau, o, \tau)
\]

\vspace{0.3em}
\textbf{Key properties:}
\begin{itemize}
    \item Deterministic straight-line trajectories (vs.\ stochastic diffusion)
    \item Generates full action chunk in parallel (no autoregressive tokens)
    \item Used by $\pi_0$, $\pi_{0.5}$, SmolVLA, and Kinetix policies
\end{itemize}

\end{frame}

%------------------------------------------------------------
\begin{frame}
    \frametitle{Appendix: Per-Environment Kinetix Breakdown ($H = 16$)}
    
    \begin{center}
        \includegraphics[width=\textwidth,height=0.80\textheight,keepaspectratio]{\kinetixfig{c16/per_environment_breakdown.pdf}}
    \end{center}
    
    \end{frame}    

\end{document}
